\documentclass[12pt]{article}
\usepackage{amsmath, amssymb}
\usepackage{graphicx}

\title{ Reproducibility and Replicability in Remote Sensing Research}
%\title{An exploratory analysis of the expressiveness of the HxC representation of pixel values for change detection}
\author{}
\date{}

\begin{document}
	
	\maketitle
	
	\section{Introduction}
	
	In recent years, ensuring the reliability of scientific research has become increasingly important. In fields such as remote sensing, where large datasets and complex models are used, it is essential that results can be trusted and verified. This report summarizes the key ideas presented in the paper \textit{``A Badging System for Reproducibility and Replicability in Remote Sensing Research''} by Frery et al., in a simplified and accessible manner.
	
	\section{What is Reproducibility and Replicability?}
	
	Before discussing the proposed system, it is important to understand two key concepts:
	
	\begin{itemize}
		\item \textbf{Reproducibility:} The ability to obtain the same results using the same data and methods as the original study.
		\item \textbf{Replicability:} The ability to achieve similar results using different data or independent methods.
	\end{itemize}
	
	In simple terms, reproducibility checks whether results can be repeated exactly, while replicability checks whether the findings are generally valid.
	
	\section{The Problem in Research}
	
	Many scientific studies, including those in remote sensing, face challenges such as:
	\begin{itemize}
		\item Lack of access to data
		\item Missing or unclear code
		\item Insufficient methodological details
	\end{itemize}
	
	These issues make it difficult for other researchers to verify results or build upon previous work.
	
	\section{Proposed Solution: A Badging System}
	
	The authors propose a \textbf{badging system} to encourage transparency and good research practices. This system assigns badges to research papers based on how well they support reproducibility and replicability.
	
	\subsection{Types of Badges}
	
	The system includes badges that indicate:
	\begin{itemize}
		\item Availability of data
		\item Availability of code
		\item Clarity of methodology
		\item Successful reproduction of results
	\end{itemize}
	
	These badges act as visual indicators that help readers quickly assess the quality and reliability of a study.
	
	\section{Benefits of the Badging System}
	
	The proposed system offers several advantages:
	\begin{itemize}
		\item Encourages researchers to share data and code
		\item Improves transparency in scientific work
		\item Helps reviewers and readers evaluate research quality
		\item Promotes trust in published results
	\end{itemize}
	
	Overall, it supports better scientific practices and more reliable findings.
	
	\section{Application in Remote Sensing}
	
	Remote sensing research often involves:
	\begin{itemize}
		\item Large datasets (e.g., satellite images)
		\item Complex algorithms
		\item Advanced statistical or machine learning methods
	\end{itemize}
	
	Because of this complexity, reproducibility is especially important. The badging system helps ensure that studies in this field can be verified and reused effectively.
	
	\section{Conclusion}
	
	The badging system proposed by Frery et al. provides a practical way to improve reproducibility and replicability in remote sensing research. By encouraging openness and transparency, it helps build trust in scientific results and supports the advancement of knowledge.
	
\end{document}